\subsection{LIMIT OF A FILTER}

DEFINITION 1. Let X be a topological space and F a filter on X. A point \(x \in X\) is said to be a limit point (or simply a limit) of F, if F is finer than the neighbourhood filter \(\mathfrak{B}(x)\) of x; F is also said to converge (or to be convergent)

to x. The point x is said to be a limit of a filter base B on X, and B is said to converge to x, if the filter whose base is B converges to x.

This definition, together with Proposition 4 of § 6, no. 3, gives the following criterion :

PROPOSITION 1. A filter base B on a topological space X converges to x if and only if every set of a fundamental system of neighbourhoods of x contains a set of B.

In accordance with the terminology introduced in § 1, no. 2 we can state Proposition 1 in the following way: B converges to x if and only if there are sets of B as near as we please to x.

If a filter \(\mathfrak{F}\) converges to \(x\) , then every filter finer than \(\mathfrak{F}\) also converges to \(x\) , by reason of Definition 1. Likewise, if the topology of \(X\) is replaced by a coarser topology, the neighbourhood filter of \(x\) is replaced by a coarser filter ( \(\S\) 2, no. 2, Proposition 3), and therefore \(\mathfrak{F}\) still converges to \(x\) in this new topology.

We can therefore say that the finer the topology, the fewer convergent filters there are in this topology. In particular, in the discrete topology, the only convergent filters are the neighbourhood filters, for these are the trivial ultrafilters on X (§ 6, no. 4).

Let \(\Phi\) be a set of filters on X, all of which converge to the same point \(x\) ; the neighbourhood filter \(\mathfrak{B}(x)\) is coarser than all the filters of \(\Phi\) , hence also coarser than the intersection \(\bigcap\) of these filters; in other words, \(\bigcap\) also converges to \(x\) .

PROPOSITION 2. A filter \(\mathfrak{F}\) on a topological space \(\mathbf{X}\) converges to a point \(x\) if and only if every ultrafilter which is finer than \(\mathfrak{F}\) converges to \(x\) .

This is an immediate consequence of the preceding remarks and Proposition 7 of § 6, no. 4.

In general a filter can have several distinct limit points; we shall revert to this question in § 8, no. 1.

